\documentclass[twoside]{article}

\usepackage{ustj}

\addbibresource{mss.bib}

\newcommand{\authorname}{Philip Monk}
\newcommand{\authorpatp}{\patp{wicdev-wisryt}}
\newcommand{\affiliation}{Tlon Corporation}

%  Make first page footer:
\fancypagestyle{firststyle}{%
\fancyhf{}% Clear header/footer
\fancyhead{}
\fancyfoot[L]{{\footnotesize
              %% We toggle between these:
              Manuscript submitted for review.\\
              % {\it Urbit Systems Technical Journal} I:1 (2024):  1–3. \\
              ~ \\
              Address author correspondence to \authorpatp.
              }}
}
%  Arrange subsequent pages:
\fancyhf{}
\fancyhead[LE]{{\urbitfont Urbit Systems Technical Journal}}
\fancyhead[RO]{I/O in Hoon}
\fancyfoot[LE,RO]{\thepage}

%%MANUSCRIPT
\title{Input and Output in Hoon}
\author{\authorname~\authorpatp \\ \affiliation}
\date{}

\begin{document}

\maketitle
\thispagestyle{firststyle}

\begin{abstract}
\end{abstract}

% We will adjust page numbering in final editing.
\pagenumbering{arabic}
\setcounter{page}{1}

\tableofcontents

\section{Introduction}

Let's talk about input and output in Urbit. I won't say much about how events and
effects are processed by the runtime; it suffices to note that the
runtime is a normal Unix program which 1) listens for various events,
2) passes them to Arvo (Urbit OS), and 3) processes the effects
produced by Arvo. Instead, I'll focus on how I/O works from the
perspective of Hoon programs, like Arvo or apps running on Arvo. But
first, let's review I/O paradigms in other languages.
\footnote{All examples are in pseudocode, since it's useful to understand
this before you're comfortable in Hoon. If you want a challenge, try translating them!}

\section{The three styles of I/O}\label{the-three-styles-of-io}

I will call these styles imperative I/O, monadic I/O, and state machines
(note that this classification refers only to the experience of using
them and the program structure they dictate). An imperative programming
language can use any of these and a functional programming language can
use any of them.

\subsection{Imperative I/O}\label{imperative-io}

Imperative languages (such as C and Python) perform I/O by calling a
function in the middle of whatever they were doing. Anywhere, in any
statement, you could have I/O.

In this example and others to follow, the task will be to read a
filename from the command line, read a \textsc{url} out of that file, and fetch
the contents at that \textsc{url}.

\begin{lstlisting}[style=listingblock]
output = fetch(stripWhitespace(readFile(cliInput())))
\end{lstlisting}

\noindent
For readability, intermediate variables would be a good idea, but
nothing stops you from nesting I/O this way.

\subsection{Monadic I/O}\label{monadic-io}

Languages that use explicit monads for I/O (such as Haskell) perform I/O
by running I/O behind a bind operation, which I'll notate
with~\texttt{←}. Thus, while you can't put
I/O~\emph{anywhere}, the general structure is a sequence of I/O
operations, such as:

\begin{lstlisting}[style=listingblock]
fileName ← cliInput()
     url ← readFile(fileName)
  output ← fetchUrl(stripWhitespace(url))
\end{lstlisting}

\noindent
Simon Peyton Jones' \href{https://www.microsoft.com/en-us/research/wp-content/uploads/2016/07/mark.pdf}{``Tackling
the Awkward Squad''} (2010) is the canonical description of monadic I/O, and does
a great job of explaining how the Haskell community decided to use this
as their primary I/O style.

\subsection{State machines}\label{state-machines}

Languages that use explicit state machines for I/O (such as Elm) perform
I/O by accepting events and producing a list of effects.

\begin{lstlisting}[style=listingblock]
event :: Initialize
       | CLIInput text
       | FileInput text
       | HTTPResponse body
effect :: CLIInput
        | ReadFile fileName
        | FetchUrl url
        | Print text
main(event) {
  switch event {
    Initialize → [CLIInput ~]
    CLIinput fileName → [ReadFile fileName]
    FileInput url → [FetchUrl stripWhitespace(url)]
    HTTPResponse body → [Print body]
  }
}
\end{lstlisting}

\noindent
The effects produced are not functions but data.

It's important to note that this example is chosen to highlight a
particular kind of I/O -- sequences of events. Further, this example has
no persistent state -- each I/O operation is a pure function of the
output of the previous one.

\section{A stateful, loopy example}\label{a-stateful-loopy-example}

Let's look at another type of program: a simple \textsc{http} server that you can
give text at the \textsc{cli}, which will respond to all \textsc{http} requests with that
text.

\newpage
\noindent
Imperative:

\begin{lstlisting}[style=listingblock]
text = "no text entered yet"
while (true)
  cli = readCli()
  if (cli != Null)
    text = cli
  request = readHttp()
  if (request != Null)
    respond(request,text)
\end{lstlisting}

\noindent
Monadic:

\begin{lstlisting}[style=listingblock]
loop("no text entered yet")
where:
loop :: text → I/O Null
  loop = input ← readInput()
  switch input
    CLIInput newText → loop(newText)
    HTTPRequest request →
      _ ← respond(request, text)
      loop(text)
\end{lstlisting}

\noindent
State machine:

\begin{lstlisting}[style=listingblock]
text = "no text entered yet"
main(text ,event) {
  switch event {
    CLIInput newText → [newText, Null]
    HTTPRequest request → [text, HTTPResponse text]
  }
}
\end{lstlisting}

\noindent
Again, note the response in the state machine is data, not a function.

I'm not trying to make a point about which type is better in general.
There's a pretty strong argument that in the first example, imperative
was the best option, followed by monadic, followed by state machine.
There's an equally strong argument that the second example is most
naturally represented by the state machine.

Look at how state is handled in the monadic example. Formally, bind
composes functions, so that every time you see
a~\texttt{←}~there's actually a new function
(closure/delimited continuation) being created and stored for when we
receive a response. All variables in scope are stored in that closure,
so we can conveniently access them later without explicitly storing
them. However, the state we preserve over time varies a lot. In this
example, it varies between
storing~\texttt{loop},~\texttt{[loop\ input]},~\texttt{[loop\ input\ newText]},
and~\texttt{[loop\ input\ request]}.

By comparison, the state machine does not store intermediate values, it
only stores its explicit state. While this is somewhat less flexible, it
puts the permanent state front-and-center.

I'll reiterate that all of these are possible to implement in purely
functional languages, and further they can all be implemented in terms
of each other. State machines are trivially implementable -- it's just a
function

\begin{lstlisting}[style=listingblock]
[state event] → [new-state (list effect)]
\end{lstlisting}

\noindent
Monadic I/O can be implemented directly (as in Haskell) or on
top of state machines by storing continuations in your state. Imperative
I/O is somewhat more difficult but should be possible through algebraic
effects.

When inventing a new system, you should consider which of these you
intend to support based on the sort of code you expect to see. For
example, a scripting language will almost always benefit from imperative
(or at least monadic) style.

\section{In Urbit}\label{in-urbit}

So much for exploring the space of I/O solutions. In terms of deployed
code on Earth, the split is about 90\% imperative, 5\% monadic, and 5\%
state machine. In Urbit, the split is about 90\% state machine and 10\%
monadic. This is what I intend to explain and defend.

Urbit's OS is called Arvo; it has several ``kernel modules'' called
vanes, and it has userspace applications. Colloquially these are called
``apps,'' but when speaking about them abstractly, it's helpful to use
their more specific
name:~\href{https://developers.urbit.org/reference/glossary/agent}{``agents''}.
Arvo, the vanes, and agents are all structured as explicit state
machines.

In addition, there
are~\href{https://developers.urbit.org/reference/glossary/thread}{``threads''},
which are structured as monadic I/O. Note there's only a superficial
resemblance to Unix threads: they are not executed in parallel and they
don't share memory, or at least not any more than agents do. They're
simply a computation that takes an argument and produces a result,
possibly after doing some I/O, structured monadically.

As you can see, all the lowest levels of Arvo are explicit state
machines, and only the highest layer supports monadic I/O. To see why
this is the case, let's list advantages and disadvantages to each type.
Note that, while I'm speaking specifically of agents and threads in the
context of Urbit, these are fundamental properties of monadic and state
machine I/O.

Agents are:

\begin{itemize}
\item
  Upgradable. Since the state is explicit, you can upgrade an agent
  in-place by supplying a new state machine and a function from the old
  state type to the new state type. This is very similar to a database
  migration.
\item
  Robust. A state machine must accept any input at any time. If it
  receives unexpected input, it must have had explicit error handling
  for that input, since it's one of the types in the switch statement.
  Thus, it's relatively easy to stay in a consistent state.
\item
  Permanent. Since it's upgradable and robust, there's no reason to stop
  its running.
\item
  Hard to write long sequences of I/O. Long sequences of I/O create many
  possible states, and you must handle any input in any state. This can
  result in complex, hard-to-follow code with code paths that are rarely
  if ever taken.
\end{itemize}

Threads are:

\begin{itemize}
\item
  Not upgradable in-place. Since the state is implicit, you can only
  cleanly upgrade by (1) killing the thread and restarting it or (2)
  letting it resolve to a defined ``upgradable state'' periodically,
  usually at the top of a ``main loop''. If you do (2), then you are in
  fact structuring it as a state machine, so the other advantages of
  threads are compromised.
\item
  Fragile. If we don't receive expected input, the default behavior is
  to be stuck waiting for input. You must remember to add any
  error-handling features you want, and it's easy to end up in an
  inconsistent state.
\item
  Impermanent. Since a thread can't be upgraded and may get stuck in
  case of unexpected input, a thread can't be expected to last forever.
  It must be ``rebooted'' from time to time.
\item
  Easy to write long sequences of I/O. Just\ldots write the I/O one line
  at a time.
\end{itemize}

Arvo proper and the vanes are permanent pieces of software and, for this
reason, state machines are more natural.

I'll note that other kernels, like the Unix kernel, use imperative I/O.
The difference is that Unix isn't permanent in itself -- it only lasts
until you reboot. Since your permanent state is stored externally, you
don't need to be upgradable or permanent. With sufficient discipline you
can write robust C code, so it makes sense to use imperative style. Even
so, a lot of C code is written in a state machine style; even if it
doesn't need to be upgradable, it's worth it for the robustness and
concurrency advantages.

However, Arvo is a single-level store -- all of its state is permanent.
This is very convenient and eliminates many long-tail bugs, but it also
means we need to code for permanence, and that's much easier when you
structure your code with explicit state.

Further, Arvo and the vanes rarely perform long sequences of I/O.
Generally, they do one thing and emit some effects, similar to the \textsc{http}
server example above. Sometimes they ``pass through'' a request. For
example, the Eyre vane passes \textsc{http} requests to userspace agents or
threads. In this case, there is a sequence of I/O, which in monadic code
would be:

loop()

where:

loop = request \texttt{←} receiveRequest() agent = lookupAgent(request)
response \texttt{←} callAgent(agent, request) \_ \texttt{←}
sendResponse(request, response) loop()

This is clear code to read, but it leaves open the question of how to
respond during the callAgent() function. If another request comes in,
can we handle it in parallel? If we need to upgrade ourselves before the
agent responds, can we make sure to properly handle the response?

The naive solution in monadic code is to say that you can only handle
one request at a time, and if an upgrade happens you just drop
outstanding requests. You can resolve both of these by introducing a
sort of main loop, which is just a way of saying ``turn it into a state
machine''.

In a state machine, the naive solution is:

state :: (map request agent) main(state, event) \{ switch event \{
\textsc{http}Request request -\textgreater{} agent = lookupAgent(request)
{[}put(state, request, agent), CallAgent agent request{]} AgentResponse
request response -\textgreater{} {[}del(state, request), \textsc{http}Respond
request response{]} \} \}

Upgrading this is trivial, since the state (map of requests to
outstanding calls to agents) is explicit. It also trivially handles
concurrent requests.

I won't claim this is easier or harder to write than the monadic
version, but it has the properties we need at this level of the system.

Userspace agents similarly are permanent entities that need to not lose
data or get stuck. For this reason they're structured as state machines.

A helpful comparison is that explicit state machines are basically Mealy
machines, the kind you might have learned about in an introductory
digital design class. This shouldn't be surprising; a digital circuit
must always be consistent because it can't be manually ``rebooted'', and
their input and output is highly formalized. By contrast, a Rube
Goldberg machine doesn't have the same constraints; which is why they're
a lot more fun!

It's notable also that explicit state machines are considerably more
declarative than imperative or monadic I/O. Not everyone thinks that's a
good thing, but anyone who defends functional programming should see
some advantages to that.

However, userspace sometimes needs to perform long sequences of I/O.
Let's look at an example where we want to fetch the front page of a site
and the first couple comments on each story. \textsc{http} is unreliable, so on
failure we retry five times.

Monadic:

topStories \texttt{←} fetch(topStoriesUrl) loop(topStories)

where:

loop :: topStories -\textgreater{} I/O Null if topStories is Null return
Null comments \texttt{←} retryLoop(5,head(topStories)) otherComments
\texttt{←} loop(tail(topStories)) append(comments,otherComments)

retryLoop :: {[}n story{]} -\textgreater{} I/O Comments comments
\texttt{←} fetchComments(story) if comments is \textsc{http}Error: if n == 0:
bail else: retryLoop(n-1, story)

State machine:

state :: Initial | FetchingTop | FetchingStory
comments stories retries | Done comments

main(state, event) \{ switch event \{ Initialize -\textgreater{} assert
state == Initial {[}FetchingTop, Fetch topStoriesUrl{]} HttpResponse
response -\textgreater{} switch state \{ FetchingTop -\textgreater{}
{[}FetchingStory Null stories 5, FetchComments head(stories){]}
FetchingStory comments stories retries -\textgreater{} if response is
\textsc{http}Error: if retries == 0: {[}Initial, Null{]} else: {[}FetchingStory
comments stories retries-1, FetchComments head(stories){]} else:
comments = append(response,comments) if tail(stories) is Null: {[}Done
comments, Null{]} else: {[}FetchingStory append(comments,response)
tail(stories) 5, FetchComments head(tail(stories)){]} \} \} \}

With practice, the state machine version of this can be written
correctly. However, (1) it's verbose, (2) the control flow jumps all
over the place, and (3) in practice it's challenging to get exactly
right. This is still a small example, and it gets much worse as the
length and complexity of the I/O sequence grows.

The monadic version, on the other hand, flows cleanly down the page. We
can factor out generic functionality like ``retry this request n times''
to make the essence of the computation clear.

Suppose you have to upgrade this. In the state machine, you can see
exactly which stage of the computation you're in and write specific code
to upgrade cleanly.

Threads, as we've discussed, can't be upgraded in-place. Let's be real
though: we're downloading comment sections, not regulating a nuclear
reactor. If we get a better version of this script, just kill the old
one and start the new one from scratch. What're a couple of extra \textsc{http}
requests?

This suggests a general division of workloads: code which needs to be
permanent, upgradable, and concurrent should generally be a state
machine. Code with long or complex sequences of I/O and definite
termination should generally be a thread.

These aren't completely mutually exclusive, but it's surprising how
often they are.

\subsection{Jael: a case study}\label{jael-a-case-study}

When you need to do something permanent and upgradable, but you also
need to do long sequences of I/O, the main trick is to factor your
problem into both an agent and a thread (or several).

As a case study, consider Jael, one of our vanes. This maintains our \textsc{pki}
state, which it downloads over \textsc{http} from an Ethereum node. Downloading
from an Ethereum node is a complex sequence that starts with fetching
its most recent block number, then fetching all the blocks we haven't
seen and scanning for transactions we care about. We also maintain a map
of block hashes to block number, and if we see that a block's parent
doesn't have the hash we expect, that means a reorganization has
occurred, so we must ``rewind'' one block at a time until we find where
the fork happened, then restart going forward.

You don't have to understand this sequence of I/O; it's sufficient to see
that it's complex and must be exactly right. We used to do this in a
state machine directly in Jael. This data is, of course, permanent and
needs to be accessible to the rest of the system, so Jael must be a
state machine.

However, we had many small mistakes in this sequence of I/O, some of
which only appeared during unusual events, such as when an \textsc{http} request
failed at the same time as a reorganization happened.

We're in a bit of a bind, though: Jael has to be a state machine. The
solution is to factor out the act of getting updates into a thread while
the main \textsc{pki} state is stored in the state machine in Jael.

This thread has a definite purpose: given the most recent block number
we already know about, fetch all the \textsc{pki} transactions since then. The
function signature is:

syncEthereum oldBlockNumber -\textgreater{} I/O {[}newBlockNumber,
newTxs{]}

Jael runs this thread every five minutes and, if it succeeds, then we
update our block number and \textsc{pki} state. If the thread fails, gets stuck,
or we need to upgrade it, we just kill the old thread and start a new
one. Again, a few extra \textsc{http} requests don't matter.

From Jael's perspective, it looks something like:

state :: State block txs outstandingThread main(state,event) \{ switch
event \{ Initialize -\textgreater{} {[}State 0 Null Null, StartTimer
(now + five minutes){]} TimerFired -\textgreater{} effects = If
outstandingThread is Null then Null else {[}Kill outstandingThread{]}
newId = newThreadId() {[}State txs newId, append(effects, StartThread
newId syncEthereum(block), StartTimer (now + five minutes)){]}
ThreadFinished Failure -\textgreater{} {[}State block txs Null, Null{]}
ThreadFinished Success newBlockNumber newTxs -\textgreater{} {[}State
newBlockNumber append(txs,newTxs) Null, Null{]} \} \}

In other words, from Jael's perspective, it's just a single I/O event
that encapsulates arbitrarily complex I/O in the thread. A single I/O
event is easy enough, and we can maintain all our permanency and
robustness guarantees, because we know the thread will have one of three
results: success, failure, or it hangs. If we set a timeout, hang
becomes failure, so there's only two possible results: success or
failure. If we handle both of those correctly, we've handled all the
possible I/O errors that could have occurred.

What about upgrading? We've already said we can handle the thread dying
for arbitrary reasons, so if an upgrade comes in, we just kill it. We
can upgrade Jael's state machine easily since its state is very simple.

\section{Recap}\label{recap}

Let's recap what we've covered:

\begin{itemize}
\tightlist
\item
  There are three styles of I/O: imperative, monadic, and explicit state
  machines.
\item
  Urbit supports explicit state machines (agents) and monadic I/O
  (threads).
\item
  Explicit state machines are preferred for permanent, robust,
  concurrent code.
\item
  Monadic I/O is preferred for short-lived code with complex I/O.
\item
  If a state machine needs to perform complex I/O, it should encapsulate
  it in a thread so that, from the perspective of the state machine,
  it's simple I/O.
\end{itemize}

Hopefully it's fairly clear why Urbit uses explicit state machines in
most situations. More than anything else, our goal is to build software
that can last forever. Explicit state machines are a crucial tool for
that.\tombstone

% \selectlanguage{USenglish}
% \printbibliography
\end{document}
